\documentclass[11pt]{article}
\usepackage[T1]{fontenc}
\usepackage{lmodern}
\usepackage{amsmath,amsthm,amssymb,booktabs,longtable}
\usepackage[hidelinks]{hyperref}
\setlength{\emergencystretch}{2em}

\newtheorem{theorem}{Theorem}[section]
\newtheorem{lemma}[theorem]{Lemma}
\newtheorem{corollary}[theorem]{Corollary}
\newcommand{\Z}{\mathbb Z}
\newcommand{\N}{\mathbb Z_{\geq 0}}
\newcommand{\one}{\mathbf 1}

\title{Twenty-seven maximum independent sets in triangular grid graphs}
\author{}
\date{2026-10-13}

\begin{document}
\maketitle

\begin{abstract}
Let \(G_3(d)\) have as vertices the nonnegative integer triples of sum
\(d\), with two vertices adjacent when their \(\ell_1\) distance is \(2\).
We prove that \(G_3(d)\) has exactly \(27\) maximum independent sets
whenever \(d\geq10\) and \(3\nmid d\). We also describe all of them:
each agrees with one of the three natural color classes away from
the three radius-three corners, and each color class has exactly nine
corner extensions. The proof uses an explicit rational triangle cover
with \(89\) weights. A coordinate reduction verifies the cover for every
allowed \(d\); saturation and a balancing argument then reduce the
classification to three ten-vertex problems. Together with Machacek's
uniqueness theorem for \(3\mid d\), the result gives Howroyd's eventual
period-three formula.
\end{abstract}

\section{Introduction and statement}\label{sec:intro}

For \(d\geq0\), put
\[
 V_d=\{(x_1,x_2,x_3)\in\N^3:x_1+x_2+x_3=d\}.
\]
The graph \(G_3(d)\) has vertex set \(V_d\), with distinct vertices \(u,v\)
adjacent if and only if \(\|u-v\|_1=2\).
Write \(\alpha_3(d)\) for its independence number and \(a_3(d)\) for the
number of independent sets of that size. Sets are counted with their
coordinates: no quotient by rotations, reflections, or graph
isomorphisms is taken. This is also the monomial graph used by
Machacek~\cite{machacek}, since, for \(u,v\in V_d\),
\[
 \deg\operatorname{lcm}(x^u,x^v)
   =d+\tfrac12\|u-v\|_1.
\]

Howroyd's conjecture, recorded in OEIS A297557 on 1 January
2018~\cite{oeis} and stated as Conjecture~1.1
in~\cite{machacek}, is
\begin{equation}\label{eq:howroyd}
 a_3(d)=
 \begin{cases}
  1,&3\mid d,\\
  27,&3\nmid d
 \end{cases}
 \qquad(d\geq9).
\end{equation}
Machacek's Theorem~2.1 establishes the divisible case.
Appendix~\ref{app:machacek} records a local correction to a deletion
statement in the cited arXiv version; it is separate from our proof of
the nondivisible case.

The independence number is already known. Machacek~\cite[Section~2]{machacek}
attributes its formula to Geramita, Gregory, and Roberts.
For \(d\geq10\) with \(3\nmid d\), it is
\(\alpha_3(d)=\binom{d+2}{2}/3\).
The same size result appears in Zabokritskiy's
Theorem~4.2~\cite{covers} and in Proposition~17 of
Kreindel, Barouch Essayag, and Zabokritskiy~\cite{codes}.
Corollary~18 of the latter paper gives uniqueness for positive degrees
divisible by \(3\), excluding degree \(6\), and cites Machacek.
These results do not count the maximum sets in the two nondivisible
residue classes. Weighted local covers and their finite-state
formulation already appear in~\cite[Section~3]{covers}.
Here a particular exact cover is used to control equality and classify
all maximum sets; the general covering method is not claimed as new.

Define the proper three-coloring
\[
 \chi(v)=v_1-v_2\pmod3,\qquad
 C_j=\{v\in V_d:\chi(v)=j\}\quad(j\in\Z/3\Z),
\]
and the body and corners
\begin{equation}\label{eq:body-corners}
 B_d=\{v\in V_d:\max_i v_i\leq d-4\},\qquad
 K_i=\{v\in V_d:v_i\geq d-3\}\quad(1\leq i\leq3).
\end{equation}
For \(d\geq10\), the \(K_i\) are disjoint and their union is
\(V_d\setminus B_d\).

\begin{theorem}\label{thm:main}
Let \(d\geq10\) and \(3\nmid d\).
Every maximum independent set of \(G_3(d)\) agrees on \(B_d\) with
exactly one color class \(C_j\).
For each \(j\), all its maximum extensions are obtained by making,
independently at the three corners, the choices in
Table~\ref{tab:corners}, in the local coordinates of
Section~\ref{sec:corners} with relative color
\(\delta_i=j-\chi(de_i)\pmod3\).
There are exactly nine extensions for each \(j\), and hence
\[
 a_3(d)=27.
\]
\end{theorem}

Thus Theorem~\ref{thm:main}, together with the cited divisible case,
gives~\eqref{eq:howroyd}. The classification is part of the counting
argument, rather than an additional assertion requiring a separate
proof.

\section{An exact triangle cover}\label{sec:cover}

Throughout Sections~\ref{sec:cover}--\ref{sec:corners}, assume
\(d\geq10\) and \(3\nmid d\), and put \(N=|V_d|=\binom{d+2}{2}\).
Every edge changes \(\chi\), so each \(C_j\) is independent.
The cyclic permutation \((v_1,v_2,v_3)\mapsto(v_2,v_3,v_1)\)
changes \(\chi\) to \(\chi-d\) modulo \(3\).
Consequently the three classes have equal size \(N/3\).

Let \(e_1,e_2,e_3\) be the standard basis vectors. For nonnegative
integer triples \(b\), define the triangles
\[
 \begin{aligned}
 U(b)&=\{b+e_1,b+e_2,b+e_3\},
       &&b_1+b_2+b_3=d-1,\\
 D(b)&=\{b+e_1+e_2,b+e_1+e_3,b+e_2+e_3\},
       &&b_1+b_2+b_3=d-2.
 \end{aligned}
\]
They are three-vertex cliques. Let \(\mathcal T_d\) be the family of
these triangles. Define
\[
 \rho(t)=
 \begin{cases}
 t,&0\leq t\leq3,\\
 4+((t-4)\bmod3),&t\geq4.
 \end{cases}
\]
The key of \(T=U(b)\) or \(D(b)\) is
\[
 \bigl(d\bmod3,\operatorname{type}(T),
       \operatorname{sort}(\rho(b_1),\rho(b_2),\rho(b_3))\bigr).
\]
Table~\ref{tab:weights} gives an integer \(W(T)\) for each of the
\(89\) keys that occur, and we set \(w(T)=W(T)/60\).

\begin{lemma}[Exact cover]\label{lem:cover}
Every triangle in \(\mathcal T_d\) has a key in
Table~\ref{tab:weights}, all its weights are nonnegative, and
\begin{equation}\label{eq:exact-cover}
 \sum_{\substack{T\in\mathcal T_d\\v\in T}}W(T)=60
 \qquad(v\in V_d).
\end{equation}
\end{lemma}

\begin{proof}
The triangles incident with \(v\) are exactly
\[
 U(v-e_i)\quad(v_i>0),\qquad
 D(v-e_i-e_j)\quad(i<j,\ v_i,v_j>0).
\]
For \(10\leq d\leq21\), \(3\nmid d\), direct integer addition using
Table~\ref{tab:weights} verifies all \(1134\) vertex equations.
All incident keys are present, all \(89\) table entries occur in this
range, and their weights are nonnegative. This finite check is
implemented in \texttt{verify\_howroyd\_certificate.py};
Section~\ref{sec:computation} specifies the computations.

If \(d>21\), some coordinate \(v_i\) is at least \(8\).
Set \(v'=v-3e_i\in V_{d-3}\).
The \(i\)-th base coordinate of any incident triangle is either \(v_i\)
or \(v_i-1\). Before and after subtraction of \(3\), that coordinate
is at least \(4\), so its \(\rho\)-value is unchanged.
The existence conditions for all incident triangles are also
unchanged, as are the other coordinates and \(d\bmod3\).
Thus the incident triangles correspond term by term, with identical
keys. Repeated reduction reaches the checked range
\(10\leq d\leq21\) without dropping below \(10\).
This proves both~\eqref{eq:exact-cover} and the absence of missing
keys at larger degrees.
\end{proof}

\begin{corollary}[Saturation]\label{cor:saturation}
The independence number is \(N/3\). If \(I\) is a maximum independent
set, then
\[
 |I\cap T|=1\qquad\text{for every \(T\in\mathcal T_d\) with \(w(T)>0\)}.
\]
\end{corollary}

\begin{proof}
For every independent set \(I\), the exact cover gives
\[
 |I|=\sum_{T\in\mathcal T_d}w(T)|I\cap T|
 \leq\sum_{T\in\mathcal T_d}w(T)
 =\frac13\sum_{v\in V_d}\sum_{T\ni v}w(T)
 =\frac N3.
\]
The color classes attain this bound. For a maximum set,
\[
 0=\frac N3-|I|
   =\sum_{T\in\mathcal T_d}w(T)(1-|I\cap T|).
\]
Every summand is nonnegative, so each triangle with positive weight
is met exactly once.
\end{proof}

\section{Propagation and constancy on the body}\label{sec:body}

The zero entries of Table~\ref{tab:weights} have a simple description.
If the actual base coordinates are sorted as \(s\leq t\leq u\), then
\begin{equation}\label{eq:zeros}
 \begin{aligned}
 W(U(b))=0
  &\ \Longleftrightarrow\
       (s,t)\in\{(0,1),(1,1),(1,2)\},\quad u\geq4,\\
 W(D(b))=0
  &\ \Longleftrightarrow\
       (s,t)\in\{(0,0),(0,2),(1,1)\},\quad u\geq4.
 \end{aligned}
\end{equation}
This equivalence is checked on the \(89\) keys. It applies to actual
coordinates because \(\rho(t)\leq2\) if and only if \(t\leq2\), and
\(\rho(t)\geq4\) if and only if \(t\geq4\).

Fix a maximum independent set \(I\), and write \(f(v)=\one_{\{v\in I\}}\).
If two positive-weight triangles share an edge, their saturation
equations show that \(f\) has the same value at their two opposite
vertices. In particular, for \(v_i\geq2\), put
\[
 F_i(v)=v+(1,1,1)-3e_i.
\]
The triangles \(U(v-e_i)\) and \(D(v-2e_i)\) share the vertices
\(v-e_i+e_j\) for \(j\ne i\); their opposite vertices are \(v\)
and \(F_i(v)\). Therefore
\begin{equation}\label{eq:propagation}
 f(v)=f(F_i(v))
\end{equation}
whenever both weights are positive. The move \(F_i\) preserves
\(\chi\).

\begin{lemma}[Balancing at a maximum coordinate]\label{lem:move}
Let \(v\in B_d\), and choose any \(i\) with \(v_i=M=\max_j v_j\).
Both \(U(v-e_i)\) and \(D(v-2e_i)\) have positive weight, and
\(F_i(v)\in B_d\).
\end{lemma}

\begin{proof}
We have \(M\geq4\), while the two other coordinates have sum at least
\(4\), since \(M\leq d-4\).
The upper triangle has base coordinates \(M-1,v_j,v_k\).
Because \(M-1\geq3\), the two small coordinates in a zero pattern
in~\eqref{eq:zeros} would have to be \(v_j,v_k\).
Their sum is at least \(4\), whereas the upper zero patterns have
small-coordinate sum at most \(3\). Thus this weight is positive.

The lower triangle has base coordinates \(M-2,v_j,v_k\).
If \(M\geq5\), the same argument applies.
If \(M=4\), the changed coordinate is \(2\).
A zero pattern not using this \(2\) is excluded by
\(v_j+v_k\geq4\).
The only lower zero pattern that can use it is \((0,2)\);
but the third original coordinate is at most \(4\), which would give
\(d\leq4+0+4=8\), a contradiction.

The decreased coordinate of \(F_i(v)\) is \(M-2\leq d-6\).
For either other coordinate \(v_j\), maximality of \(M\) gives
\(2v_j\leq M+v_j\leq d\), hence
\[
 v_j+1\leq\lfloor d/2\rfloor+1\leq d-4
 \qquad(d\geq10).
\]
All new coordinates are nonnegative. Thus \(F_i(v)\) remains in \(B_d\).
\end{proof}

\begin{lemma}\label{lem:constant}
For each \(j\in\Z/3\Z\), \(f\) is constant on \(B_d\cap C_j\).
\end{lemma}

\begin{proof}
Write \(d=3k+r\), with \(r\in\{1,2\}\); here \(k\geq3\).
Whenever the maximum coordinate is \(M\geq k+2\), perform a move
from Lemma~\ref{lem:move}. For the nonnegative integer potential
\(\Phi(v)=v_1^2+v_2^2+v_3^2\), the change is
\[
 \Phi(F_i(v))-\Phi(v)=2d-6M+6\leq2r-6<0.
\]
The process therefore terminates, preserving the body, color, and
value of \(f\).

At termination all coordinates are at most \(k+1\).
For \(r=2\), the only possibilities are the three permutations of
\((k,k+1,k+1)\); they have distinct colors.
For \(r=1\), the possibilities are the permutations of
\((k+1,k,k)\) and \((k-1,k+1,k+1)\).
There is one point of each color in each of these two families.
Matching points are joined by a maximum-coordinate move, as in
\[
 (k+1,k,k)\longmapsto(k-1,k+1,k+1).
\]
These moves are allowed by Lemma~\ref{lem:move}; no decrease of the
potential is needed for this final connection.
Consequently all starting points of a given color have the same
value of \(f\).
\end{proof}

\begin{corollary}[One color on the body]\label{cor:body}
There is exactly one \(j\in\Z/3\Z\) such that
\[
 I\cap B_d=C_j\cap B_d.
\]
\end{corollary}

\begin{proof}
For \(r=1\), the triangle \(U(k,k,k)\) lies in \(B_d\) and has
positive weight. For \(r=2\), use \(U(k+1,k,k)\).
Indeed \(k\geq3\), their base coordinates exclude the zero
patterns~\eqref{eq:zeros}, and their vertex coordinates are at most
\(k+2\leq d-4\).
Each triangle has one vertex of each color. The three binary
constants from Lemma~\ref{lem:constant} sum to \(1\) by saturation.
Exactly one is therefore \(1\), and all three color classes occur
in the body.
\end{proof}

\section{The corner configurations and the count}\label{sec:corners}

The graph distance from \(v\) to \(de_i\) is \(d-v_i\).
Thus each \(K_i\) is a radius-three corner with ten vertices.
For vertices in distinct corners their graph distance is at least
\(d-6\geq4\), so there are no edges between corners.
Use the following coordinate maps:
\begin{equation}\label{eq:corner-maps}
 \begin{aligned}
 \psi_1(p,q)&=(d-p-q,p,q),\\
 \psi_2(p,q)&=(q,d-p-q,p),\\
 \psi_3(p,q)&=(p,q,d-p-q).
 \end{aligned}
\end{equation}
In each case \(p,q\geq0\) and \(p+q\leq3\) describe the corner;
adjacency is given by the six differences
\(\pm(1,0),\pm(0,1),\pm(1,-1)\).
The only body vertices adjacent to the corner are on \(p+q=4\).
That whole layer lies in \(B_d\): its coordinates are \(d-4,p,q\)
with \(p,q\leq4\leq d-4\).

Put \(c_i=\chi(de_i)\), so \((c_1,c_2,c_3)=(d,-d,0)\) modulo \(3\).
Direct substitution in~\eqref{eq:corner-maps} gives
\[
 \chi(\psi_i(p,q))=c_i+p-q\pmod3.
\]
For a fixed body color \(j\), the selected vertices of the outer
layer therefore satisfy
\[
 p-q\equiv\delta_i\pmod3,\qquad \delta_i=j-c_i.
\]
In particular, relative color \(0\) means that the background color
contains the corner vertex.

\begin{lemma}[Corner enumeration]\label{lem:corners}
Fix \(\delta\in\Z/3\Z\) and the vertices of \(p+q=4\) with
\(p-q\equiv\delta\pmod3\).
Among independent sets in \(p,q\geq0,\ p+q\leq3\) having no
neighbor among those fixed vertices, all sets of maximum size are
exactly those in Table~\ref{tab:corners}.
\end{lemma}

\begin{proof}
This is a finite enumeration on ten vertices.
For \(\delta=0,1,2\), respectively, deleting vertices adjacent to the
fixed outer layer leaves \(8,7,7\) available vertices.
Enumerate all subsets of the available vertices, discard those
containing one of the six adjacency differences, and retain those
of largest size. The resulting sets are precisely the seven rows
of Table~\ref{tab:corners}.
The original certificate verifier implements this enumeration
directly from the corner graph and its fixed outer layer.
An independent implementation also enumerates all \(2^{10}\)
subsets before applying the same independence and compatibility
conditions.
\end{proof}

\begin{table}[htbp]
\centering
\small
\begin{tabular}{ccl}
\toprule
Relative color \(\delta\)&Maximum size&All maximum corner sets\\
\midrule
0&4&\(\{(0,0),(0,3),(1,1),(3,0)\}\)\\
\midrule
1&3&\(\{(0,0),(0,2),(2,0)\}\)\\
 & &\(\{(0,0),(0,2),(2,1)\}\)\\
 & &\(\{(0,2),(1,0),(2,1)\}\)\\
\midrule
2&3&\(\{(0,0),(0,2),(2,0)\}\)\\
 & &\(\{(0,0),(1,2),(2,0)\}\)\\
 & &\(\{(0,1),(1,2),(2,0)\}\)\\
\bottomrule
\end{tabular}
\caption{Complete corner choices compatible with the fixed outer layer.
Coordinates are \((p,q)\) as in~\eqref{eq:corner-maps}.}
\label{tab:corners}
\end{table}

\begin{proof}[Proof of Theorem~\ref{thm:main}]
Corollary~\ref{cor:body} determines a unique body color \(j\).
Since \(3\nmid d\), the three corner colors \(d,-d,0\) are distinct,
and hence the relative colors \(\delta_i\) are \(0,1,2\) in some order.
The background color has \(4,3,3\) vertices in these respective
corners, so it has \(N/3-10\) vertices in the body.

Every maximum independent set has size \(N/3\).
By Lemma~\ref{lem:corners}, the sum of its three corner sizes is at
most \(4+3+3=10\); equality in the total size forces each corner
to attain its local maximum. There are consequently
\(1\cdot3\cdot3=9\) choices for the fixed body color.
Conversely, every such choice is compatible with the body;
distinct corners have no edges between them. It is an independent
set of size \(N/3\), hence is maximum.
Different choices within one background give different sets, and
different backgrounds are already distinct on the body.
There are therefore \(3\cdot9=27\) maximum independent sets, with
exactly the stated classification.
\end{proof}

\section{Computation and reproducibility}\label{sec:computation}

The final certificate consists solely of the \(89\) nonnegative
integer numerators in Table~\ref{tab:weights}, also supplied as
\texttt{howroyd\_cover\_weights.csv}; every denominator is \(60\).
Floating-point linear programming was used during the search for
weights. Neither its output tolerances nor a numerical optimality
claim are assumptions of the proof. The original verifier uses
integer arithmetic and finite-set operations, reads the CSV, and
does not import the search code.

With the two Python programs and CSV in the same directory, run
the following commands using Python~3.10 or later and its standard
library:
\begingroup
\footnotesize
\begin{verbatim}
python verify_howroyd_certificate.py
python referee_checks.py structural --csv howroyd_cover_weights.csv
python referee_checks.py dp --max-degree 20
\end{verbatim}
\endgroup
The structural command uses its default maximum degree \(60\).
The explicit \texttt{--csv} argument avoids dependence on the
independent checker's original \texttt{submission/} directory layout.

The original verifier checks the \(89\) entries, all \(1134\)
base vertex equations, the exact zero-weight patterns, and the
corner enumeration. Its further diagnostics check \(34\,800\)
vertex equations at degrees \(31,32,49,50,100,101,202\) and
\(24\,684\) body moves at nondivisible degrees from \(10\) through \(60\).
These additional diagnostics are not needed for the
all-degree reduction in Lemma~\ref{lem:cover}.

The separately written program \texttt{referee\_checks.py} does
not import the original verifier.
Its structural mode reconstructs adjacency and finds triangles
as graph cliques. For \(10\leq d\leq60\), \(3\nmid d\), it checks
\(25\,704\) vertex equations, \(47\,821\) triangles, and \(25\,296\)
maximum-coordinate moves, including every choice when maxima tie.
It also checks the body color components and the corner
configurations; for the \(16\) nondivisible degrees from \(10\)
through \(32\), it constructs and checks the \(27\) distinct
maximum-size extensions. The structural mode still uses the same
weight certificate.

Its full-graph dynamic program is a separate cross-check.
It processes rows of lengths \(1,2,\ldots,d+1\), recording the
largest size and its number of realizations for each independent
row mask. A short-row mask \(s\) and the following long-row mask
\(t\) are compatible exactly when
\[
 s\mathbin{\&}\bigl(t\mathbin{|}(t\mathbin{\gg}1)\bigr)=0,
\]
where the operations are bitwise. Each vertex set has a unique
row representation. This computation does not use the weights,
body classification, or a prescribed answer \(27\).
It was run for every \(d=0,\ldots,20\); in particular it gives
\[
 a_3(7)=28,\quad a_3(8)=32,\quad
 a_3(9)=a_3(12)=a_3(15)=a_3(18)=1
\]
and \(a_3(d)=27\) for
\(d\in\{10,11,13,14,16,17,19,20\}\).
The original verifier, structural checks, and dynamic program
were rerun locally and reproduced the recorded results.
These are finite checks. The assertions for every allowed degree
use the coordinate reduction and balancing proofs above.

\section{Provenance and verification status}\label{sec:provenance}

AI systems contributed to the mathematical argument, certificate,
verification programs, and manuscript preparation.
The argument underwent AI-assisted review, and the verification
programs were rerun locally with the results described above.
The \(27\)-count theorem for all \(d\geq10\) with \(3\nmid d\)
has also been formalized in Lean~4.19.0 using mathlib v4.19.0.
The source, fixed dependency revisions, and build instructions are
provided in the repository's \texttt{lean/} directory.
The final theorem depends only on Lean's standard axioms
\texttt{propext}, \texttt{Classical.choice}, and \texttt{Quot.sound};
it uses no unproved placeholders, custom axioms, or native evaluation.
The cited divisible-by-three branch is not part of this formalization.
No human expert review or checking by a separate external proof kernel
has been completed.
The literature search completed for this draft did not locate an
earlier proof of the full \(27\)-set statement; that search was not
exhaustive and does not establish priority.
The references below identify the versions examined. In particular,
Machacek's published journal text was not obtained for comparison
with the cited arXiv version.

\appendix
\section{A boundary-deletion correction in the divisible case}
\label{app:machacek}

This note concerns only the cited proof of the known divisible
case, not Theorem~\ref{thm:main}.
On page~5 of the arXiv version of~\cite{machacek}, the proof of
Theorem~2.1 takes \(d=3k\), identifies a forced boundary independent
set \(J\) of size \(3k=d\), and writes
\(G'''=G_3(d)\setminus J\cong G_3(d-3)\).
With ordinary vertex deletion, this is not literally correct.
The entire boundary
\[
 \partial V_d=\{v\in V_d:\min_i v_i=0\}
\]
has \(3d=9k\) vertices for \(d>0\). Deleting only \(J\) therefore
leaves \(6k\) more vertices than \(G_3(d-3)\).

The required correction is local to that step.
The preceding argument on each boundary side \(X_i\) gives
\(|I\cap X_i|=k+1\) and forces \(J\cap X_i\subseteq I\cap X_i\)
with \(|J\cap X_i|=k+1\). Hence
\[
 I\cap\partial V_d=J.
\]
Delete the whole boundary and put
\[
 H=G_3(d)[V_d\setminus\partial V_d].
\]
Subtracting \((1,1,1)\) gives \(H\cong G_3(d-3)\).
Equation~(2.3) of~\cite{machacek} then gives
\[
 |I\setminus J|
   =\alpha_3(d)-3k
   =\alpha_3(d-3)
   =\alpha(H).
\]
Thus \(I\setminus J\) is a maximum independent set of \(H\), and
the stated induction continues: its inductive step has \(k>3\),
so \(d-3\geq9\).
This is an explicit correction to the wording of the examined
version, not a quotation of a correction made by its author or
a new attribution of the uniqueness theorem.

\section{The complete integer certificate}\label{app:weights}

In Table~\ref{tab:weights}, \(r=d\bmod3\), \(\tau\) is the triangle
type, and \(a\leq b\leq c\) are the sorted \(\rho\)-values of its
base coordinates, not the unmodified coordinates. The final column
is \(W=60w\). The table contains all \(89\) rows of
\nolinkurl{howroyd_cover_weights.csv}.

\small
\begin{longtable}{cccccr}
\caption{Integer numerators of the exact triangle cover.}
\label{tab:weights}\\
\toprule
\(r\)&\(\tau\)&\(a\)&\(b\)&\(c\)&\(W\)\\
\midrule
\endfirsthead
\multicolumn{6}{c}{Table \thetable\ (continued)}\\
\toprule
\(r\)&\(\tau\)&\(a\)&\(b\)&\(c\)&\(W\)\\
\midrule
\endhead
\midrule
\multicolumn{6}{r}{Continued on next page}\\
\endfoot
\bottomrule
\endlastfoot
1&D&0&0&5&0\\
1&D&0&1&4&30\\
1&D&0&2&6&0\\
1&D&0&3&5&5\\
1&D&0&4&4&10\\
1&D&0&5&6&3\\
1&D&1&1&6&0\\
1&D&1&2&5&22\\
1&D&1&3&4&13\\
1&D&1&4&6&15\\
1&D&1&5&5&16\\
1&D&2&2&4&3\\
1&D&2&3&3&24\\
1&D&2&3&6&17\\
1&D&2&4&5&12\\
1&D&2&6&6&10\\
1&D&3&3&5&3\\
1&D&3&4&4&24\\
1&D&3&5&6&10\\
1&D&4&4&6&17\\
1&D&4&5&5&17\\
1&D&5&6&6&17\\
1&U&0&0&6&60\\
1&U&0&1&5&0\\
1&U&0&2&4&30\\
1&U&0&3&6&30\\
1&U&0&4&5&25\\
1&U&0&6&6&32\\
1&U&1&1&4&0\\
1&U&1&2&6&0\\
1&U&1&3&5&3\\
1&U&1&4&4&4\\
1&U&1&5&6&3\\
1&U&2&2&5&16\\
1&U&2&3&4&3\\
1&U&2&4&6&8\\
1&U&2&5&5&6\\
1&U&3&3&3&3\\
1&U&3&3&6&17\\
1&U&3&4&5&3\\
1&U&3&6&6&17\\
1&U&4&4&4&3\\
1&U&4&5&6&3\\
1&U&5&5&5&3\\
1&U&6&6&6&3\\
2&D&0&0&6&0\\
2&D&0&1&5&30\\
2&D&0&2&4&0\\
2&D&0&3&6&3\\
2&D&0&4&5&3\\
2&D&0&6&6&6\\
2&D&1&1&4&0\\
2&D&1&2&6&24\\
2&D&1&3&5&15\\
2&D&1&4&4&6\\
2&D&1&5&6&8\\
2&D&2&2&5&3\\
2&D&2&3&4&24\\
2&D&2&4&6&17\\
2&D&2&5&5&3\\
2&D&3&3&3&3\\
2&D&3&3&6&3\\
2&D&3&4&5&10\\
2&D&3&6&6&3\\
2&D&4&4&4&3\\
2&D&4&5&6&3\\
2&D&5&5&5&3\\
2&D&6&6&6&3\\
2&U&0&0&4&60\\
2&U&0&1&6&0\\
2&U&0&2&5&30\\
2&U&0&3&4&30\\
2&U&0&4&6&27\\
2&U&0&5&5&30\\
2&U&1&1&5&0\\
2&U&1&2&4&0\\
2&U&1&3&6&3\\
2&U&1&4&5&9\\
2&U&1&6&6&7\\
2&U&2&2&6&12\\
2&U&2&3&5&3\\
2&U&2&4&4&3\\
2&U&2&5&6&17\\
2&U&3&3&4&3\\
2&U&3&4&6&10\\
2&U&3&5&5&3\\
2&U&4&4&5&17\\
2&U&4&6&6&17\\
2&U&5&5&6&17\\
\end{longtable}
\normalsize

\begin{thebibliography}{9}

\bibitem{oeis}
OEIS Foundation Inc.,
\emph{The On-Line Encyclopedia of Integer Sequences},
entry A297557.
Howroyd's conjecture dated 1 January 2018 and Machacek's comment
dated 1 May 2021.
\url{https://oeis.org/A297557/internal}.
Consulted 13 October 2026.

\bibitem{machacek}
John Machacek,
\emph{Unique maximum independent sets in graphs on monomials of a
fixed degree},
arXiv:2010.11112v2, 14 March 2021.
\url{https://arxiv.org/abs/2010.11112v2}.

\bibitem{covers}
Aryeh Lev Zabokritskiy (Yohananov),
\emph{Independent Sets in Multiset Profile Graphs via Weighted
Local Covers},
arXiv:2607.13733v3, 4 August 2026.
\url{https://arxiv.org/abs/2607.13733v3}.

\bibitem{codes}
Avraham Kreindel, Isaac Barouch Essayag, and
Aryeh Lev Zabokritskiy (Yohananov),
\emph{Multiset Deletion Codes: Cyclic Constructions, Bounds, and
Exact Results},
arXiv:2601.05636v2, 13 July 2026.
\url{https://arxiv.org/abs/2601.05636v2}.

\end{thebibliography}
\end{document}
